% Discussion for item 4.
% scripts/make_latex.py will NEVER overwrite this file -- edit it freely.

The output matches the form the problem sheet asks for: the child announced
itself, the parent announced itself, and the parent then printed the message the
child had sent through the pipe.

\texttt{pipe(fd)} produces one buffer inside the kernel with two handles onto it,
\texttt{fd[0]} to read and \texttt{fd[1]} to write, and \texttt{fork()} copies the
whole descriptor table --- so immediately afterwards \emph{four} descriptors refer
to that one pipe. The child closes \texttt{fd[0]} and the parent closes
\texttt{fd[1]}, leaving exactly one writer and one reader. That closing is not
housekeeping: the kernel signals end-of-file to a reader only when \emph{every}
write end has been closed, so a parent still holding its own copy of
\texttt{fd[1]} would wait forever for an EOF that its own descriptor prevents,
and the program would simply hang.

The pipe does guarantee that the third line comes after the first, since the
parent's \texttt{read()} cannot return until the child has written. It guarantees
nothing about the second line: the child's announcement and the parent's are
printed by two runnable processes that are not synchronised with each other, so
their order is left to the scheduler, exactly as in item 1.1. Our program inserts
a brief \texttt{usleep()} in the parent purely so the output matches the expected
listing, which hides the race rather than resolving it --- a fixed delay is a
guess about timing, and a wrong one under load. The pipe already provides real
synchronisation, and using it, for instance by having the parent print only after
its \texttt{read()} returns, would order the output without any sleep at all.
